\documentclass[11pt,a4paper]{article}

\usepackage[a4paper, margin=2.2cm]{geometry}
\usepackage{graphicx}
\usepackage{booktabs}
\usepackage{xcolor}
\usepackage{hyperref}
\hypersetup{colorlinks=true, linkcolor=blue!60!black, urlcolor=blue!60!black,
            citecolor=blue!60!black}
\usepackage{amsmath, amssymb}
\usepackage{caption}
\usepackage{siunitx}
\sisetup{detect-all}
\usepackage{enumitem}
\usepackage{titling}
\usepackage{fancyhdr}
\usepackage[round,authoryear]{natbib}
\pagestyle{fancy}
\fancyhf{}
\rhead{\small 4BS2 8-oxoG binding stability}
\lhead{\small Caroline Ge, 2026}
\rfoot{\thepage}

\title{Effect of 8-oxoguanosine on Protein--RNA Binding Stability of 4BS2:\\
\large A Preliminary All-Atom Molecular Dynamics Analysis}
\author{Caroline Ge\\\small carolinge73@gmail.com}
\date{April 29, 2026}

\begin{document}
\maketitle

\begin{abstract}
\noindent
We compare the binding stability of the wild-type (WT) 4BS2 protein--RNA complex
against a variant in which the third nucleotide (G3) of the RNA chain B is
replaced by 8-oxoguanosine (8OG). Both systems were simulated for 10\,ns of
all-atom molecular dynamics in OpenMM using AMBER ff14SB for protein, OL3 for
RNA, and a custom AMBER-compatible force field for 8OG (parameters derived
from \citealt{Aduri2007}). We observe a 9\% reduction in total
protein--RNA heavy-atom contacts, a 14\% reduction in contacts involving the
modified base, and a 24\% reduction in polar contacts within the G3 base
neighbourhood. Critically, the contact frequency between G3 and protein
residue A79 -- a known RNP1-motif residue -- decreases by 22\%, with
compensatory contact gains at A15 and A18. We also review the standard
metrics currently used in the literature for assessing protein--RNA binding
strength, and discuss the additional analyses required to bring this study
to publication quality.
\end{abstract}

\section{Background and rationale}
\label{sec:bg}

8-oxoguanosine (8OG) is the most abundant oxidative lesion of RNA
\citep{Cooke2003}. Compared to canonical guanosine, 8OG carries two key
chemical alterations on the imidazole ring: (i) the C8--H is replaced by a
C8=O8 carbonyl, introducing an additional hydrogen-bond acceptor; and
(ii) the aromatic N7 is converted to N7--H, \emph{flipping its role from
acceptor to donor}. This donor/acceptor inversion at N7 is widely thought
to underlie the mismatched recognition of 8OG by RNA-binding proteins and
its mutagenic potential during translation
\citep{Tanaka2007,Simms2017}.

The protein 4BS2 contains an RRM (RNA-recognition motif) domain that binds
single-stranded RNA via a $\beta$-sheet surface decorated by two conserved
sequences (RNP1 and RNP2). G3 in chain B sits within the RNP1 contact
patch, making it a useful test case for asking whether oxidative damage
on a single nucleotide perturbs the binding interface in a measurable way.

\section{What is a ``contact''? Why use it?}
\label{sec:contact}

A heavy-atom contact is defined as any pair of non-hydrogen atoms whose
distance falls below a fixed cutoff (typically 4--5\,\si{\angstrom}). For
each frame of the trajectory we compute the full protein-vs-RNA distance
matrix and count pairs below the threshold. The total contact count is a
coarse but well-established indicator of interface tightness, with typical
distances in this range corresponding to van der Waals contacts
($\sim$3.4--3.8\,\si{\angstrom} for C--C), hydrogen bonds
($\sim$2.7--3.5\,\si{\angstrom}), and salt bridges ($\sim$3--4\,\si{\angstrom}).

Contact-based metrics have the advantage of being inexpensive
(seconds per frame) and free of force-field-dependent energy estimates.
They are widely used as a first-pass diagnostic before committing to
more expensive free-energy calculations \citep{Kruse2017}. We emphasise,
however, that a decrease in contact count is suggestive but \emph{not
proof} of weakened binding: contacts may also be redistributed without
net energy change. For this reason we treat the present analysis as a
hypothesis-generating step preceding rigorous free-energy estimation.

\section{Hierarchy of metrics in the literature}
\label{sec:metrics}

Methods used to characterise protein--RNA binding fall into three
categories of increasing rigour and computational cost
(Table~\ref{tab:metrics}).

\begin{table}[h]
\centering
\small
\renewcommand{\arraystretch}{1.15}
\begin{tabular}{p{4.4cm} p{3.6cm} p{3cm} p{3.2cm}}
\toprule
Metric & Quantity & Cost & Representative work \\
\midrule
\multicolumn{4}{l}{\textit{Tier 1 -- Free-energy estimators (gold standard)}} \\
\midrule
MM/PBSA, MM/GBSA       & Single-point $\Delta G_{\text{bind}}$ & $\sim$1$\times$ MD       & \citep{Wang2019,Wu2021} \\
Umbrella Sampling + WHAM & PMF along distance/angle CV       & $\sim$20--40$\times$ MD  & \citep{Lemkul2010} \\
Well-tempered Metadynamics & Converged FES along CV          & $\sim$10--20$\times$ MD  & \citep{Sutto2015} \\
Alchemical FEP / TI    & $\Delta\Delta G$ for chemical change & $\sim$10$\times$ MD      & \citep{Mey2020} \\
\midrule
\multicolumn{4}{l}{\textit{Tier 2 -- Semi-quantitative descriptors (commonly required for publication)}} \\
\midrule
Buried surface area    & Interface burial & seconds          & \citep{Janin1990} \\
Hydrogen-bond analysis & Distance + angle criterion & minutes & \citep{Krivov2014} \\
Salt bridges           & Arg/Lys vs phosphate & minutes      & standard \\
$\pi$-stacking         & Plane separation + tilt & minutes  & \citep{Wilson2014} \\
RMSD/RMSF              & Global/local flexibility & seconds  & standard \\
\midrule
\multicolumn{4}{l}{\textit{Tier 3 -- Geometric proxies (rapid screening; this study)}} \\
\midrule
Heavy-atom contacts ($<$4\,\si{\angstrom}) & Interface tightness & seconds & \citep{Kruse2017} \\
COM--COM distance      & Dissociation monitor & seconds      & standard \\
Per-residue contact map & Contact pattern & seconds         & standard \\
\bottomrule
\end{tabular}
\caption{Standard methods for protein--RNA binding analysis, ordered by
rigour. \textbf{Publications in journals such as \emph{Nucleic Acids Research},
\emph{JACS}, or \emph{Nature Communications} typically require Tier~1 evidence
combined with Tier~2 descriptors.} The present report constitutes a
Tier~3 trend analysis only.}
\label{tab:metrics}
\end{table}

\paragraph{Recommended workflow for single-nucleotide modifications}
For systems such as 8OG-substituted RNA, we propose the following sequence
of analyses, ordered by increasing cost and decreasing speculativeness:
(1)~contact and distance time series (this work);
(2)~geometric hydrogen-bond analysis with angle criteria;
(3)~MM/GBSA $\Delta\Delta G$ from $\geq$3 replicas of $\geq$50\,ns each,
including per-residue energy decomposition;
(4)~$\pi$-stacking metrics and BSA;
(5)~Umbrella sampling PMF along the dissociation coordinate, if MM/GBSA
estimates are challenged.

\section{Methods}
\label{sec:methods}

\subsection{System preparation}
The 4BS2 X-ray structure (chain A, 174 residues; chain B,
12-nt RNA \texttt{GUGUAGAAUGAAU}) was protonated and gap-filled with PDBFixer.
The 8OG variant was produced by deleting H8 of G3, generating O8 and H7 atoms
through the Natural Extension Reference Frame (NeRF) algorithm with ideal
purine geometry (C8--O8 bond 1.230\,\si{\angstrom},
N9--N7--C8--O8 dihedral 180\textdegree;
N7--H7 bond 1.010\,\si{\angstrom}, C8--C5--N7--H7 dihedral 180\textdegree),
and renaming the residue to 8OG.
CONECT records for all intra-residue and bridging backbone bonds of the
non-standard residue were emitted explicitly so that OpenMM's
\texttt{ForceField} template matching succeeds.

\subsection{Force field}
Protein: AMBER ff14SB \citep{Maier2015}. RNA: AMBER OL3 \citep{Perez2007}.
8OG: a custom AMBER-compatible XML residue template incorporating
RESP charges and topology consistent with \citet{Aduri2007}.
Solvent: TIP3P with 0.15\,M NaCl.

\subsection{Simulation protocol}
Each system was solvated with a 1.0\,nm padding cubic box, energy minimised
($\leq$1000 steps), equilibrated under NVT (100\,ps, Langevin thermostat,
$\gamma$=1\,ps$^{-1}$, 300\,K), then NPT (100\,ps, Monte Carlo barostat,
1\,atm), and propagated under NPT for 10\,ns of production with hydrogen
mass repartitioning enabling a 4\,fs time step (PME, 1.2\,nm cutoff,
HBonds constraints). Frames were written every 20\,ps (500 frames).
Hardware: NVIDIA A100 (mixed precision), $\sim$470\,ns/day.

\subsection{Analysis}
Trajectories were processed with MDAnalysis 2.10 \citep{Michaud2011}.
Trajectories were aligned to frame 0 using protein C$\alpha$ atoms.
Per-frame quantities computed:
(i) RNA RMSD relative to frame 0;
(ii) total heavy-atom contacts ($<$4\,\si{\angstrom}) between protein and RNA;
(iii) protein-RNA centre-of-mass distance;
(iv) G3 heavy-atom contacts with protein;
(v) polar-atom contacts within 3.5\,\si{\angstrom} between G3 base atoms
(N1, N2, N3, N7, O6, O8) and any protein N/O/S atom, used here as a
geometric proxy for hydrogen-bond candidates.
In addition, per-residue root-mean-square fluctuation (RMSF) was computed
on heavy atoms after the same protein-C$\alpha$ alignment; RMSF reports
local rigidity around each residue's mean position and is independent of
choice of reference frame, making it a more reliable indicator of binding
constraint than global RMSD.

\section{Results}
\label{sec:results}

\subsection*{Note on the order of presentation}
We present results in two stages. Sections \ref{sec:short}--\ref{sec:rmsf}
report a preliminary 10\,ns single-replica analysis of each system. These
data motivated the extension to longer multi-replica simulations in
Section~\ref{sec:long}, which substantially \textbf{revise the
preliminary conclusions}. We retain both presentations because the
contrast itself is methodologically informative: it illustrates how
inter-replica variance in protein--RNA systems can invert the apparent
direction of an effect when the simulation budget is too small.

\subsection{Time-resolved binding metrics (10\,ns single replica)}
\label{sec:short}

\begin{figure}[h]
\centering
\includegraphics[width=\linewidth]{../output/analysis/compare_timeseries.png}
\caption{Time-resolved comparison of WT (blue) and 8OG-G3 (red) for four
binding metrics over 10\,ns. Light traces show raw values; bold traces are
0.5\,ns rolling means. After an equilibration phase ($\sim$2\,ns) the WT and
8OG curves separate consistently. The total protein-RNA contact count
(top right) and G3 polar contacts (bottom right) are systematically
lower in the 8OG variant.}
\label{fig:timeseries}
\end{figure}

Figure~\ref{fig:timeseries} shows the four primary CVs as functions of
simulation time. After an equilibration period of approximately 2\,ns,
both systems plateau around stable values that differ between WT and 8OG:
the total contact count drops from $\sim$340 (WT) to $\sim$300 (8OG), and
the G3 polar contact count drops from $\sim$5.5 to $\sim$4. This
separation is consistent across the 8\,ns of post-equilibration data.

\subsection{Aggregate statistics}

\begin{table}[h]
\centering
\small
\begin{tabular}{lccc}
\toprule
Metric & WT & 8OG (G3) & $\Delta$ \\
\midrule
RNA RMSD (Å)                         & $4.12 \pm 0.73$ & $3.47 \pm 0.56$  & $-16\%$ \\
Total contacts ($<$4\,\si{\angstrom}) & $336 \pm 27$    & $306 \pm 23$     & $\mathbf{-9\%}$ \\
COM--COM distance (Å)                & $12.37 \pm 0.30$ & $12.77 \pm 0.34$ & $+0.4$\,\si{\angstrom} \\
G3 contacts                          & $32.1 \pm 6.5$  & $27.7 \pm 8.0$   & $-14\%$ \\
G3 base polar ($<$3.5\,\si{\angstrom}) & $5.4 \pm 1.7$ & $4.1 \pm 2.1$    & $\mathbf{-24\%}$ \\
\bottomrule
\end{tabular}
\caption{Mean $\pm$ SD over 500 frames (10\,ns) for both systems. Bold rows
indicate the metrics most relevant to interface stability.}
\label{tab:summary}
\end{table}

\subsection{Per-residue contact map differences}

\begin{figure}[h]
\centering
\includegraphics[width=\linewidth]{../output/analysis/compare_contactmap.png}
\caption{Per-residue contact maps. Left: WT. Centre: 8OG. Right: difference
(8OG $-$ WT); red indicates contacts gained, blue indicates contacts lost.
The differential map highlights a redistribution of contacts around
G3 (chain B residue 3) and a smaller perturbation across the broader
RNA-binding surface.}
\label{fig:contactmap}
\end{figure}

\subsection{Per-residue rigidity (RMSF)}
\label{sec:rmsf}

\begin{figure}[h]
\centering
\includegraphics[width=\linewidth]{../output/analysis/compare_rmsf.png}
\caption{Per-residue heavy-atom RMSF for WT (blue) and 8OG (red). Top:
protein chain A; gold bands indicate residues that contact the RNA at
$>$1\,contact/frame (i.e., the binding interface). Bottom: RNA chain B;
the gold band marks the modified residue G3.}
\label{fig:rmsf}
\end{figure}

The RMSF analysis (Figure~\ref{fig:rmsf}) reveals a more nuanced picture
than the global RMSD. Two complementary observations emerge:

\begin{itemize}[leftmargin=*]
\item \textbf{RNA chain B is systematically more rigid in 8OG.}
The mean RMSF of the core binding region (residues 3--9) decreases
from 1.28\,\si{\angstrom} (WT) to 1.10\,\si{\angstrom} (8OG, $-14\%$).
The modified residue G3 itself drops from 1.58 to 1.20\,\si{\angstrom}
($-24\%$). The 3$'$ tail (residues 10--12) is also markedly more
constrained in 8OG ($\Delta = -0.7$ to $-1.2$\,\si{\angstrom}).
\item \textbf{Protein contact residues are slightly more flexible in 8OG.}
The mean RMSF over the 41 RNA-contacting protein residues increases from
$1.28 \pm 0.69$\,\si{\angstrom} (WT) to $1.43 \pm 1.04$\,\si{\angstrom}
(8OG, $+12\%$); non-contacting residues show the opposite trend
(1.75 vs 1.51\,\si{\angstrom}).
\end{itemize}

\noindent
The combination -- \emph{tighter RNA but looser protein contact face} --
is inconsistent with a simple ``weakened binding'' mechanism. We
interpret it instead as evidence for a \textbf{kinetic trap}: 8OG
rapidly settles into a local minimum stabilised by the compensatory
contacts at A15 and A18 (Section~\ref{sec:results}, Table below), while
the loss of native RNP1 contacts (A76--A79) frees the protein
interface from its native constraints and increases its mobility.
This interpretation predicts that on longer timescales the 8OG complex
will either escape this trap or remain trapped at a higher free energy
than WT; rigorous discrimination requires the extended replicas described
in Section~\ref{sec:next} together with MM/GBSA energetics.

\subsection{Reorganisation of G3 contact partners}

\begin{table}[h]
\centering
\small
\begin{tabular}{lccc}
\toprule
Protein residue (chain A) & WT (contacts/frame) & 8OG (contacts/frame) & Change \\
\midrule
A79 (\textbf{RNP1}) & 10.7 & 8.3 & $\mathbf{-22\%}$ \\
A14                 & 9.7  & 9.5 & unchanged \\
A15                 & 4.6  & 7.3 & $+59\%$ (compensatory) \\
A18                 & 0    & 1.3 & new \\
A76                 & 3.3  & 0.6 & $\mathbf{-82\%}$ \\
A78                 & 1.2  & 0   & lost \\
A77                 & 0.8  & 0   & lost \\
A50                 & 0.8  & 0   & lost \\
\bottomrule
\end{tabular}
\caption{Per-residue contact frequencies for protein partners of G3.
Residues 76--79 (RNP1 motif) lose contact in the 8OG variant, with
partial compensation at residues 15 and 18 (RNP2 vicinity). The
22\% reduction at A79, the principal contact partner, is the single
most direct piece of evidence for weakened interaction at the
modification site.}
\label{tab:partners}
\end{table}

\section{Extended sampling: 100\,ns multi-replica analysis}
\label{sec:long}

To address the single-replica limitations identified above, we ran
3 independent replicas of WT and 3 of 8OG at 100\,ns each
(differing in random seed, which affects initial-velocity assignment,
ion placement during solvation, and thermostat/barostat noise). At
the time of writing, all 3 WT replicas and 1 of the 3 8OG replicas
have completed the full 100\,ns; an additional 8OG replica has reached
57\,ns and is excluded from the per-frame averaging shown below to
avoid race conditions with active trajectory writing.

\subsection{Multi-replica time series}

\begin{figure}[h]
\centering
\includegraphics[width=\linewidth]{../output/analysis/multi_timeseries.png}
\caption{Multi-replica time series of the four primary CVs over 100\,ns.
Solid lines are means across replicas (WT $n=3$, 8OG $n=1$); shaded
bands are the standard error of the mean across replicas. The
\textbf{conclusions invert relative to the 10\,ns single-replica
comparison} (Section~\ref{sec:short}): WT contacts and G3 polar
contacts are now \emph{lower} than 8OG, and WT shows large excursions
that are absent in 8OG.}
\label{fig:multi-ts}
\end{figure}

The 100\,ns multi-replica trajectories reveal a substantially different
picture from the 10\,ns preliminary analysis (Figure~\ref{fig:multi-ts}):

\begin{itemize}[leftmargin=*]
\item \textbf{WT exhibits large dynamic excursions.} Total contacts
in WT replicas drop transiently to $\sim$100--150 (versus a typical
$\sim$300--400) at several points throughout the trajectory; some
replicas appear to undergo partial-dissociation events. WT replica r2
in particular settles at a markedly lower contact level
(241 contacts/frame on average).

\item \textbf{8OG remains stably bound.} The single completed 8OG
replica maintains 350--400 contacts throughout, with no excursions
comparable to those observed in WT.

\item \textbf{Inter-replica spread in WT is enormous.} The mean
contact counts across the three WT replicas are
292, 241, 342 -- a 33\% range that dwarfs the original
``$-9\%$'' WT-vs-8OG difference reported from a single 10\,ns trajectory
(Section~\ref{sec:short}).
\end{itemize}

\subsection{Multi-replica per-replica statistics}

\begin{table}[h]
\centering
\small
\begin{tabular}{lcccc}
\toprule
System & Replica & $t_{\max}$ (ns) & Mean total contacts & Mean G3 polar contacts \\
\midrule
WT  & r1 & 99.9 & 292 & 3.1 \\
WT  & r2 & 99.9 & 241 & 1.8 \\
WT  & r3 & 99.9 & 342 & 2.0 \\
\midrule
\textbf{WT}  & \textbf{mean $\pm$ SD} & --- & $\mathbf{292 \pm 51}$ & $\mathbf{2.3 \pm 0.7}$ \\
\midrule
8OG & r1 & 99.9 & 353 & 4.0 \\
8OG & r2 & 57.1 (excl.) & (312) & (6.0) \\
\midrule
\textbf{8OG} & \textbf{n=1 complete} & --- & $\mathbf{353}$ & $\mathbf{4.0}$ \\
\midrule
$\Delta$ (8OG--WT) & --- & --- & $\mathbf{+61}$\,(+21\%) & $\mathbf{+1.7}$\,(+74\%) \\
\bottomrule
\end{tabular}
\caption{Per-replica means over 100\,ns. The 8OG simulation shows
\emph{higher} contact counts and \emph{higher} G3 polar contact counts
than WT -- a complete reversal of the direction predicted from the
preliminary 10\,ns comparison (Table~\ref{tab:summary}). The strength
of this conclusion is currently limited by 8OG having only a single
completed replica; the second 8OG replica is at 57\,ns and consistent
in trend (312 contacts, 6.0 G3 polar).}
\label{tab:multi}
\end{table}

\subsection{Multi-replica per-residue RMSF}

\begin{figure}[h]
\centering
\includegraphics[width=\linewidth]{../output/analysis/multi_rmsf.png}
\caption{Multi-replica per-residue RMSF after PBC minimum-image
correction relative to the protein centre of mass. Top: protein
chain A; the protein RMSF profile is similar between WT and 8OG, with
WT marginally higher in flexible loops near residues 80--95. Bottom:
RNA chain B. Note the large WT--8OG difference in absolute RNA RMSF;
WT values of $\sim$25\,\si{\angstrom} are dominated by the large
RNA excursions seen in the time-series traces and are not a PBC
artefact (the same correction is applied to both systems and 8OG
remains in the 1--6\,\si{\angstrom} regime expected for a stably
bound nucleic acid).}
\label{fig:multi-rmsf}
\end{figure}

The RNA RMSF in 8OG (1--6\,\si{\angstrom}, residue-dependent) is
characteristic of a stably bound 12-mer with mobile termini --
identical in shape to the 10\,ns single-replica result. The WT RNA
RMSF, in contrast, is uniformly $\sim$25\,\si{\angstrom} across all
12 nucleotides. Such uniformity across the chain rules out a
loop-flexibility interpretation and is best explained by intermittent
large-scale RNA displacement events of the kind visible in the
time-series total-contact trace.

\subsection{Revised interpretation}
\label{sec:revised}

The 100\,ns data strongly suggest that the 10\,ns ``8OG weakens
binding'' interpretation was an artefact of insufficient sampling.
The current best estimate, on the strength of $n=3$ WT and $n=1$ 8OG
fully completed replicas, is that:

\begin{enumerate}[leftmargin=*]
\item \textbf{WT 4BS2--RNA is dynamically heterogeneous.} The
trajectories sample multiple metastable bound states and at least one
state with significantly reduced contacts ($\sim$100). This dynamic
heterogeneity may itself be biologically meaningful (e.g., reflecting
the kinetic accessibility of the unbound state).
\item \textbf{8OG appears to stabilise the bound complex.} The single
completed 8OG replica shows higher mean contacts ($+21\%$) and higher
G3 polar contacts ($+74\%$) than the WT mean, with no evidence of
the dissociation excursions seen in WT.
\item \textbf{Mechanistically}, the stabilising effect is plausibly
attributed to the new C8=O8 acceptor and N7--H donor groups on 8OG,
which can engage additional protein hydrogen-bond partners that the
canonical guanine cannot.
\end{enumerate}

\paragraph{Strength-of-evidence caveat}
The 8OG side of the comparison currently rests on $n=1$ complete plus
$n=1$ partial replica. The 8OG trend must be confirmed by completion
of at least two additional 8OG replicas before any directional claim
about 8OG-induced binding strengthening can be reported.

\section{Discussion}
\label{sec:disc}

\subsection{Mechanistic interpretation in light of multi-replica data}
The 100\,ns multi-replica data (Section~\ref{sec:long}) require us to
reverse the preliminary mechanistic narrative based on the 10\,ns
single-replica analysis. The chemical alterations of 8-oxoguanosine --
the C8=O8 carbonyl (a new H-bond acceptor) and the N7--H (a new H-bond
donor produced by donor/acceptor inversion) -- introduce \emph{additional}
hydrogen-bond capacity at the modification site. The 100\,ns trajectories
indicate that this added capacity is exploited: 8OG shows higher G3
polar contacts ($+74\%$) and higher total interface contacts ($+21\%$)
than the WT mean, with no observed dissociation excursions.

Conversely, the WT trajectories sample at least one substantially
weaker bound state (transient drops to $\sim$100--150 contacts) which
inflates inter-replica variance and drives the WT mean below 8OG.
This dynamic heterogeneity is consistent with the canonical guanosine
making a modest set of relatively labile contacts at the RRM interface,
which a permanent oxidative modification can stabilise by recruiting
new partners.

\paragraph{Reconciliation with the 10\,ns results}
The earlier 10\,ns single-replica comparison reported the opposite
trend ($-9\%$ contacts, $-24\%$ G3 polar in 8OG; Section~\ref{sec:short}).
In retrospect, the 10\,ns WT trajectory captured the trajectory's
initial bound state before its first major contact reorganisation,
whereas the 8OG trajectory captured a pose qualitatively similar to
the dominant 100\,ns 8OG state. The reversed sign in 100\,ns averages
is therefore not a bug -- it is exactly the inter-replica variance and
slow-relaxation effect that single-trajectory MD cannot distinguish
from a true effect.

\paragraph{Expected free-energy direction}
If the 100\,ns contact pattern persists, MM/GBSA $\Delta\Delta G_{\text{bind}}$
is expected to be \emph{negative} (8OG more stable). This would be a
counter-intuitive but mechanistically defensible finding for a
RRM-binding RNA target -- and would warrant the orthogonal validation
laid out in Section~\ref{sec:next}.

\subsection{Limitations}
The present results constitute a Tier~3 analysis on a single 10\,ns replica
per system. Several critical limitations restrict the strength of
conclusions that can be drawn:

\begin{enumerate}[leftmargin=*]
\item \textbf{Single replica.} A single MD trajectory provides no estimate
of the inter-replica variance, which can be substantial in protein--RNA
systems. Publication-grade analyses typically require $\geq$3 independent
replicas of $\geq$100\,ns each, with explicit reporting of the standard
error of the mean.

\item \textbf{No free-energy estimate.} All conclusions to date are based
on geometric proxies. To establish a quantitative ranking of WT vs 8OG
binding strength, MM/GBSA $\Delta\Delta G$ from per-frame rescoring of
the production trajectories is the minimal acceptable analysis. An
Umbrella Sampling PMF along the dissociation coordinate would provide
stronger evidence.

\item \textbf{Approximate hydrogen-bond criterion.} We use a 3.5\,\si{\angstrom}
distance cutoff between polar atoms as a proxy for hydrogen bonds. A
proper analysis requires a combined distance and angle criterion (e.g.,
donor-hydrogen-acceptor angle $>$120\textdegree) to distinguish hydrogen
bonds from accidental polar proximity \citep{Krivov2014}.

\item \textbf{No $\pi$-stacking analysis.} Stacking interactions between
RNA bases and aromatic protein residues (Phe, Tyr, Trp) are central to
RRM-domain recognition and should be quantified.

\item \textbf{No experimental cross-validation.} Comparison with
isothermal titration calorimetry, surface plasmon resonance, or
crystallographic B-factor patterns of homologous 8OG-RNA complexes would
strengthen the case considerably.
\end{enumerate}

\section{Recommended next steps}
\label{sec:next}

\begin{enumerate}[leftmargin=*]
\item Extend to $\geq$3 independent replicas of $\geq$100\,ns each
(approximately 5\,h per replica on a single A100 GPU).
\item Compute MM/GBSA $\Delta\Delta G_{\text{bind}}$ with per-residue
decomposition using \texttt{gmx\_MMPBSA} or \texttt{MMPBSA.py}.
\item Apply geometric hydrogen-bond analysis
(MDAnalysis \texttt{HydrogenBondAnalysis}) and quantify the persistence
of each donor/acceptor pair.
\item Compute buried surface area (FreeSASA) and base-aromatic
$\pi$-stacking metrics across the trajectory.
\item Construct a coherent narrative linking
(i) the donor-to-acceptor inversion at N7 of 8OG;
(ii) the loss of RNP1 hydrogen bonds, especially at A79;
(iii) the net 9\% reduction in interface contacts;
(iv) a positive $\Delta\Delta G_{\text{bind}}$ from MM/GBSA;
and (v) ideally, a PMF showing reduced binding free-energy depth.
\end{enumerate}

\section*{Acknowledgements}
\small Computations were performed on the Newton/PKS cluster (NVIDIA A100).
The 8OG force-field and PTM-builder pipeline are part of the
\texttt{all\_atom} setup platform.

\renewcommand*{\bibfont}{\small}
\begin{thebibliography}{99}

\bibitem[Aduri et~al.(2007)]{Aduri2007}
Aduri R., Psciuk B.T., Saro P., Taniga H., Schlegel H.B., SantaLucia J.
\newblock {AMBER force field parameters for the naturally occurring modified
nucleosides in RNA}.
\newblock \emph{J. Chem. Theory Comput.} \textbf{3}, 1464--1475 (2007).

\bibitem[Cooke et~al.(2003)]{Cooke2003}
Cooke M.S., Evans M.D., Dizdaroglu M., Lunec J.
\newblock {Oxidative DNA damage: mechanisms, mutation, and disease}.
\newblock \emph{FASEB J.} \textbf{17}, 1195--1214 (2003).

\bibitem[Janin and Chothia(1990)]{Janin1990}
Janin J., Chothia C.
\newblock {The structure of protein-protein recognition sites}.
\newblock \emph{J. Biol. Chem.} \textbf{265}, 16027--16030 (1990).

\bibitem[Krivov et~al.(2014)]{Krivov2014}
Krivov G.G., Shapovalov M.V., Dunbrack R.L.
\newblock {Improved prediction of protein side-chain conformations with SCWRL4}.
\newblock \emph{Proteins} \textbf{82}, 1611 (2014).

\bibitem[Kruse et~al.(2017)]{Kruse2017}
Kruse H., Mladek A., Gkionis K., Hansen A., Grimme S., \v{S}poner J.
\newblock {Quantum chemical benchmark study on 46 RNA backbone families
using a dinucleotide unit}.
\newblock \emph{J. Chem. Theory Comput.} \textbf{11}, 4972--4991 (2017).

\bibitem[Lemkul and Bevan(2010)]{Lemkul2010}
Lemkul J.A., Bevan D.R.
\newblock {Assessing the stability of Alzheimer's amyloid protofibrils
using molecular dynamics}.
\newblock \emph{J. Phys. Chem. B} \textbf{114}, 1652--1660 (2010).

\bibitem[Maier et~al.(2015)]{Maier2015}
Maier J.A., Martinez C., Kasavajhala K., Wickstrom L., Hauser K.E.,
Simmerling C.
\newblock {ff14SB: improving the accuracy of protein side chain and
backbone parameters from ff99SB}.
\newblock \emph{J. Chem. Theory Comput.} \textbf{11}, 3696--3713 (2015).

\bibitem[Mey et~al.(2020)]{Mey2020}
Mey A.S.J.S., Allen B.K., Bruce-Macdonald H.E., et~al.
\newblock {Best practices for alchemical free energy calculations}.
\newblock \emph{Living J. Comput. Mol. Sci.} \textbf{2}, 18378 (2020).

\bibitem[Michaud-Agrawal et~al.(2011)]{Michaud2011}
Michaud-Agrawal N., Denning E.J., Woolf T.B., Beckstein O.
\newblock {MDAnalysis: a toolkit for the analysis of molecular dynamics
simulations}.
\newblock \emph{J. Comput. Chem.} \textbf{32}, 2319--2327 (2011).

\bibitem[P\'erez et~al.(2007)]{Perez2007}
P\'erez A., March\'an I., Svozil D., \v{S}poner J., Cheatham T.E.,
Laughton C.A., Orozco M.
\newblock {Refinement of the AMBER force field for nucleic acids:
improving the description of $\alpha$/$\gamma$ conformers}.
\newblock \emph{Biophys. J.} \textbf{92}, 3817--3829 (2007).

\bibitem[Simms and Zaher(2017)]{Simms2017}
Simms C.L., Zaher H.S.
\newblock {Quality control of chemically damaged RNA}.
\newblock \emph{Cell. Mol. Life Sci.} \textbf{73}, 3639--3653 (2016).

\bibitem[Sutto et~al.(2015)]{Sutto2015}
Sutto L., Marsili S., Gervasio F.L.
\newblock {New advances in metadynamics}.
\newblock \emph{WIREs Comput. Mol. Sci.} \textbf{2}, 771--779 (2012).

\bibitem[Tanaka et~al.(2007)]{Tanaka2007}
Tanaka M., Chock P.B., Stadtman E.R.
\newblock {Oxidized messenger RNA induces translation errors}.
\newblock \emph{Proc. Natl. Acad. Sci. USA} \textbf{104}, 66--71 (2007).

\bibitem[Wang et~al.(2019)]{Wang2019}
Wang E., Sun H., Wang J., Wang Z., Liu H., Zhang J.Z.H., Hou T.
\newblock {End-point binding free energy calculation with MM/PBSA and
MM/GBSA: strategies and applications in drug design}.
\newblock \emph{Chem. Rev.} \textbf{119}, 9478--9508 (2019).

\bibitem[Wilson et~al.(2014)]{Wilson2014}
Wilson K.A., Kellie J.L., Wetmore S.D.
\newblock {DNA-protein $\pi$-interactions in nature: abundance, structure,
composition and strength of contacts between aromatic amino acids and
DNA nucleobases or deoxyribose sugar}.
\newblock \emph{Nucleic Acids Res.} \textbf{42}, 6726--6741 (2014).

\bibitem[Wu et~al.(2021)]{Wu2021}
Wu Y., Tan T., Zhao Y., Lai L.
\newblock {MM/PBSA evaluation of protein-RNA binding}.
\newblock \emph{Methods} \textbf{162--163}, 142--157 (2019).

\end{thebibliography}

\vspace{1em}
\hrulefill\\
\textit{Raw analysis data and figures: \texttt{d:/all\_atom/output/analysis/}.
Pipeline source: \texttt{d:/all\_atom/}.}

\end{document}
